\documentclass[10pt,a4paper]{amsart}
\usepackage[margin=19mm]{geometry}
\usepackage{amsmath,amssymb,mathtools}
\usepackage[scr=rsfs,cal=euler]{mathalfa}
\usepackage{xfrac}
\usepackage[unicode,hidelinks,bookmarks=false]{hyperref}
\newcommand{\ie}{i.e.\ }
\newcommand{\cf}{cf.\ }
\newcommand{\resp}{respectively\ }
\newcommand{\Subsection}{\subsection}
\providecommand{\nobreakdash}{\nobreak\hbox{-}}
\setlength{\emergencystretch}{2em}
\multlinegap0pt
\usepackage{amssymb}
\usepackage{mathtools}
\usepackage{xcolor}
\usepackage{tikz}
\usepackage[shortlabels]{enumitem}
\newtheorem{theorem}{Theorem}
\newcommand{\Exc}{\mathcal{E}}
\renewcommand{\phi}{\varphi}
\title{Homomorphism Counts Quadratic Residues}
\author{Priyabrata Mandal$^*$, Sonu Kumar, Deep Bhattacharjee}
\date{Source: arXiv:2502.14266v4 (CC BY 4.0)}
\begin{document}
\maketitle

\noindent\emph{Source and licence.} Homomorphism Counts Quadratic Residues. Version arXiv:2502.14266v4 (CC BY 4.0), distributed by arXiv under the Creative Commons Attribution 4.0 International licence, http://creativecommons.org/licenses/by/4.0/, as recorded in the arXiv licence metadata for this exact version and shown on the abstract page https://arxiv.org/abs/2502.14266v4. Full licence notice: \texttt{licenses/arxiv-2502.14266-CC-BY-4.0.txt}.\par\smallskip

\noindent\textit{Editorial scope.} Counted unit: the article's theorem thm:div (source0.tex lines 375--379). The article's complete original proof of it is delivered with it (source0.tex lines 381--426, one \texttt{proof} environment of 46 lines). No mathematics is written, repaired or re-derived: the statement and the proof are the article's own lines, in the article's own notation, and no retained line is substituted or re-flowed.  No further source-local context is needed: every cross-reference of every retained line resolves inside the two delivered spans.  Nothing was omitted from any retained span, so the package carries no omission record.  No retained line contains a citation, so the wrapper carries no bibliography and no bibliography entry.  No retained line contains a figure environment, an image-inclusion command or drawing code, so no image is delivered and the package declares no asset dependency.  Editorial formatting only, in addition: an AMS \texttt{amsart} preamble that restates the article's own declarations and macros used by the retained lines, namely the environments \texttt{theorem} on separate counters, together with the standard packages those lines need.  The retained environments print in this document as Theorem 1 in order of appearance, whereas the article prints the same items with the numbering of its own full document.  The complete native source of the article as distributed in the arXiv e-print for this exact version is bundled unchanged as \texttt{sources/173570/source0.tex}.
\begin{theorem}[Divisibility criterion]\label{thm:div}
  Let $n>2$ and $n\mid m$.  Then
  $2^{\omega(n)}\mid\phi(n)$
  if and only if $n\notin\Exc$.
\end{theorem}

\begin{proof}
Track the $2$-adic valuation $v_2(\phi(n))$ against $\omega(n)$.

\smallskip\noindent
\textbf{Case~I ($n$ odd).}\;
Write $n=p_1^{r_1}\cdots p_k^{r_k}$, $\omega(n)=k$.
Since each $p_i$ is odd, $2\mid(p_i-1)$, giving
\[
  v_2(\phi(n))=v_2\Bigl(\prod_{i=1}^k p_i^{r_i-1}(p_i-1)\Bigr)
  \;\ge\; k=\omega(n).
\]
Hence $2^{\omega(n)}\mid\phi(n)$.

\smallskip\noindent
\textbf{Case~II ($n=2^s q$, $s\ge2$, $q$ odd).}\;
$\omega(n)=\omega(q)+1$.
\[
  \phi(n)=\phi(2^s)\phi(q)=2^{s-1}\phi(q).
\]
Since $s\ge2$, $v_2(2^{s-1})\ge1$, and Case~I gives
$v_2(\phi(q))\ge\omega(q)$.  Thus $v_2(\phi(n))\ge1+\omega(q)=\omega(n)$.

\smallskip\noindent
\textbf{Case~III ($n=2\alpha$, $\alpha$ odd).}\;
$\omega(n)=\omega(\alpha)+1$ and $\phi(n)=\phi(\alpha)$.
Write $\alpha=p_1^{r_1}\cdots p_k^{r_k}$.  Then
$v_2(\phi(\alpha))=\sum_{i=1}^k v_2(p_i^{r_i-1}(p_i-1))=\sum_{i=1}^k v_2(p_i-1)$.

Since each $p_i$ is an odd prime,
$v_2(p_i-1)\ge1$; equality holds iff $p_i\equiv3\pmod4$.
Thus $v_2(\phi(n))=\sum_i v_2(p_i-1)$.

We need $v_2(\phi(n))\ge k+1=\omega(n)$.  This holds iff
$\sum_i v_2(p_i-1)\ge k+1$, i.e.\ at least one $v_2(p_i-1)\ge2$,
i.e.\ at least one $p_i\equiv1\pmod4$.

Conversely, if every $p_i\equiv3\pmod4$ then $v_2(p_i-1)=1$ for each
$i$, giving $v_2(\phi(n))=k<k+1=\omega(n)$, so
$2^{\omega(n)}\nmid\phi(n)$.
Precisely this condition defines $n\in\Exc$.\footnote{%
  The condition is symmetric: $n=2\alpha$ fails divisibility for
  \emph{every} product $\alpha$ of primes $\equiv3\pmod4$, regardless
  of multiplicities.  Thus $18=2\cdot3^2$, $54=2\cdot3^3$, and
  $98=2\cdot7^2$ are all exceptions, as the table confirms.}
\qedhere
\end{proof}

\end{document}
